\documentclass[../main.tex]{subfiles}

\begin{document}

\section{Chapter 4: Chart Types for Amounts, Distributions, Proportions, Relationships, and Trends}
\label{sec:chapter4-chart-types}

\subsection{Techniques / Principles Applied}

\begin{itemize}[leftmargin=*]
\item Matching chart types to amounts, distributions, proportions,
relationships, and trends.
\item Preferring readable 2D Cartesian layouts for most dashboard comparisons.
\end{itemize}

\subsection{How Applied in the Dashboard}

\subsubsection{Matching Chart Types to Analytical Questions}

\dashboardfigure{0.56}{fig-page4-cash-only-reasons.png}
{Why consumers stick to cash-only transactions.}
{fig:cash-only-reasons}

\textbf{Amounts.} \figref{fig:cash-only-reasons}{Cash-only reasons} uses a
standard bar chart to compare the weighted size of cash-only barriers. This is
appropriate because the task is simple magnitude ranking, and bar length on a
common baseline makes the differences easy to read.

\dashboardfigure{0.62}{fig-page1-population-pyramid-age-gender.png}
{The weighted population pyramid by age and gender.}
{fig:population-pyramid-age-gender}

\textbf{Distributions.} \figref{fig:population-pyramid-age-gender}{Population pyramid by age and gender}
and \figref{fig:socioeconomic-matrix}{Socioeconomic matrix} show how the
population is distributed across age, gender, income, and education. These
views are suitable because they reveal whether the sample is balanced or
concentrated in specific groups before later policy comparisons begin.

\textbf{Proportions.} \figref{fig:national-inflows}{National inflow digitization}
uses a 100\% stacked bar layout to compare payment composition within each
inflow type. This chart works well when the goal is to compare how cash,
direct-to-account, and other channels share the same total.

\textbf{Relationships.} \figref{fig:mobile-ownership-profile}{Mobile ownership profile}
uses a small-multiples bar design to compare ownership patterns across
location, gender, income, and phone type. The structure avoids overplotting
while still showing how the variables relate to one another.

\textbf{Trends.} \figref{fig:save-borrow-trajectories}{Save vs. borrow trajectories}
uses line overlays to show how saving and borrowing rates move across ordered
income groups and income runway categories. The chart is effective for showing
direction of change even though the categories are ordinal rather than
continuous.

\subsubsection{Preferring Readable 2D Cartesian Layouts}

Across \figref{fig:cash-only-reasons}{Cash-only reasons},
\figref{fig:population-pyramid-age-gender}{Population pyramid by age and gender},
\figref{fig:national-inflows}{National inflow digitization}, and
\figref{fig:save-borrow-trajectories}{Save vs. borrow trajectories}, the
dashboard relies mainly on a 2D Cartesian coordinate system because
perpendicular axes support clearer comparison than radial or 3D alternatives
for this type of policy dashboard.

\noindent\textit{Statistical uncertainty is not explicitly visualized, so the
charts are easy to read but can overstate precision, especially for filtered
subgroup views.} This matters most when the reader narrows the dashboard to a
small cohort, because visually clear differences may look firmer than the
survey design can guarantee without confidence intervals or sample-size cues.

\subsection{Notes / Adjustments}

\figref{fig:national-inflows}{National inflow digitization} emphasizes
composition rather than exact comparison of every internal segment, so only the
first segment in each row has a shared baseline.
\figref{fig:save-borrow-trajectories}{Save vs. borrow trajectories} also
simplifies ordinal data into continuous-looking trajectories, which improves
readability but should be interpreted as directional rather than strictly
continuous change.

\end{document}
